% Generated by task tex-catalogue from dossiers/code-spine.md, section A (catalogue).
% Snippets are the dossier's, extracted by it from `git show b435631:<file>`.
\section{The spine, declaration by declaration}\label{sec:gen-cat-spine}

Every public declaration of the modules in scope, in file order, with its class: \textbf{load-bearing} (the statement of a master theorem unfolds to it; \enquote{named form} marks those reached only through \texttt{piop\_\allowbreak{}rbrKnowledgeSoundness}'s named extractor and state function, section B), \textbf{interface} (a phase author must meet or use it), \textbf{helper} (used only inside proofs; those a master statement unfolds to are marked), \textbf{test fixture}. Snippets are extracted from \texttt{git show b435631:<file>} by \texttt{p\allowbreak{}r\allowbreak{}o\allowbreak{}b\allowbreak{}e\allowbreak{}s\allowbreak{}/\allowbreak{}c\allowbreak{}o\allowbreak{}d\allowbreak{}e\allowbreak{}-{}\allowbreak{}s\allowbreak{}p\allowbreak{}i\allowbreak{}n\allowbreak{}e\allowbreak{}/\allowbreak{}c\allowbreak{}a\allowbreak{}t\allowbreak{}a\allowbreak{}l\allowbreak{}o\allowbreak{}g\allowbreak{}u\allowbreak{}e\allowbreak{}.\allowbreak{}p\allowbreak{}y}: structures and inductives in full, definitions with their data body and without proof fields, theorems and instances as their statement. Helpers that are pure proof lemmas (\texttt{p\allowbreak{}a\allowbreak{}s\allowbreak{}s\allowbreak{}T\allowbreak{}h\allowbreak{}r\allowbreak{}o\allowbreak{}u\allowbreak{}g\allowbreak{}h\allowbreak{}\_\allowbreak{}m\allowbreak{}a\allowbreak{}t\allowbreak{}e\allowbreak{}r\allowbreak{}i\allowbreak{}a\allowbreak{}l\allowbreak{}i\allowbreak{}z\allowbreak{}e\allowbreak{}O\allowbreak{}u\allowbreak{}t\allowbreak{}p\allowbreak{}u\allowbreak{}t}, \texttt{p\allowbreak{}a\allowbreak{}s\allowbreak{}s\allowbreak{}T\allowbreak{}h\allowbreak{}r\allowbreak{}o\allowbreak{}u\allowbreak{}g\allowbreak{}h\allowbreak{}V\allowbreak{}e\allowbreak{}r\allowbreak{}i\allowbreak{}f\allowbreak{}i\allowbreak{}e\allowbreak{}r\allowbreak{}\_\allowbreak{}t\allowbreak{}o\allowbreak{}V\allowbreak{}e\allowbreak{}r\allowbreak{}i\allowbreak{}f\allowbreak{}i\allowbreak{}e\allowbreak{}r\allowbreak{}\_\allowbreak{}r\allowbreak{}u\allowbreak{}n}, \texttt{passThroughPure}, \texttt{passThroughProver}, \texttt{passThroughVerifier}, \texttt{passThroughExtractor}, \texttt{passThroughStateFunction}, \texttt{passThrough\_\allowbreak{}rbr}, \texttt{sendSpec}, \texttt{sendProver}, \texttt{sendEmbedding}, \texttt{sendVerifier}, \texttt{sendOracle\_\allowbreak{}outputPure}, \texttt{send\_\allowbreak{}materializeOutput}, \texttt{sendVerifier\_\allowbreak{}toVerifier\_\allowbreak{}run}, \texttt{sendVerifierPure}, \texttt{sendOracle\_\allowbreak{}complete}, \texttt{sendWitMid}, \texttt{sendOracle\_\allowbreak{}rbr}, \texttt{Refinement.\allowbreak{}id}, \texttt{Refinement.\allowbreak{}comp}, \texttt{noOracle\_\allowbreak{}eq}, and the seventeen lemmas of \texttt{KnowledgeAppend} other than its two named declarations) are named here and not listed one by one; the send-oracle ones are counted in section B where the named statement reaches them.

Library objects (all at the pins ArkLib \texttt{dca90385}, CompPoly \texttt{3468b38c}, VCVio \texttt{f9dc47d9}) are introduced where first used in sections C and D; the catalogue names them in the \enquote{depends on} column.

\par\noindent{\small\textit{Cross-references.} A reference such as \enquote{(D.5)} or \enquote{section B} in this part is to a section of the dossier \file{dossiers/code-spine.md}, whose codes the headings keep in parentheses: A (catalogue) is \cref{sec:gen-cat-spine}; B (audit surface) is \cref{sec:gen-surface-spine}; C (the statements) is \cref{sec:gen-statements-spine}; D (non-vacuity by probe), G (what was not verified) and the appendix (the probes) is \cref{sec:gen-probes-spine}; E (conformance) is \cref{sec:gen-conf-spine}; F (findings) is \cref{sec:gen-findings-spine}.}\par

Order of this section: the five spine modules (\texttt{Instance}, \texttt{Seams}, \texttt{Phase}, \texttt{Compose}, \texttt{Toy}), then the generic modules under \texttt{ToArkLib/}, then \texttt{Field.\allowbreak{}lean} where the spine unfolds to it. Every load-bearing and interface declaration is shown with its snippet, its source at \texttt{b435631}, and its meaning and dependencies as the dossier states them; the helpers and test fixtures of each file are listed in a table at the end of the file's subsection. The class is given at the right of each entry.

\subsection{\texttt{LeanerVM/Protocol/Spine/Instance.lean}}\label{sec:gen-cat-spine-Instance}

\par\medskip\noindent\textbf{\lean{Side}}\hfill\textsf{\small [load-bearing]}\par\nopagebreak

\begin{leancode}
inductive Side
  | push
  | pull
  deriving DecidableEq, Repr
\end{leancode}

\src{LeanerVM/Protocol/Spine/Instance.lean:51-54 at b435631}

Which side of the bus a tuple is on: provided (\texttt{push}) or consumed (\texttt{pull}). \emph{Depends on:} —.

\par\medskip\noindent\textbf{\lean{Shape}}\hfill\textsf{\small [load-bearing]}\par\nopagebreak

\begin{leancode}
structure Shape where
  ntab : ℕ
  τ : Fin ntab → ℕ
  width : Fin ntab → ℕ
\end{leancode}

\src{LeanerVM/Protocol/Spine/Instance.lean:57-63 at b435631}

The tables: how many, each one's log-height (the announced sizes) and width. \emph{Depends on:} —.

\par\medskip\noindent\textbf{\lean{Shape.ColumnId}}\hfill\textsf{\small [load-bearing]}\par\nopagebreak

\begin{leancode}
abbrev Shape.ColumnId (S : Shape) : Type := Σ j : Fin S.ntab, Fin (S.width j)
\end{leancode}

\src{LeanerVM/Protocol/Spine/Instance.lean:66-66 at b435631}

A column: a table and a column index in it. \emph{Depends on:} \texttt{Shape}.

\par\medskip\noindent\textbf{\lean{Layout}}\hfill\textsf{\small [load-bearing]}\par\nopagebreak

\begin{leancode}
structure Layout (μ : ℕ) (ι : Type) (κ : ι → ℕ) where
  read : Column μ → (c : ι) → Column (κ c)
  extend : (c : ι) → Vector E (κ c) → Vector E μ
  read_eval : ∀ (q : Column μ) (c : ι) (z : Vector E (κ c)),
    CMlPolynomialEval.eval₂Mle (read q c).values (algebraMap K E) z =
      CMlPolynomialEval.eval₂Mle q.values (algebraMap K E) (extend c z)
\end{leancode}

\src{LeanerVM/Protocol/Spine/Instance.lean:73-81 at b435631}

How a column is read off the stack (\texttt{read}) and how a point of the column lifts to a point of the stack (\texttt{extend}), with the one law: the column's extension at \texttt{z} is the stack's at \texttt{extend c z}. A reading law: it does not say the columns are disjoint slices (C.2, D.4). \emph{Depends on:} \texttt{Column}, CompPoly \texttt{eval₂Mle}.

\par\medskip\noindent\textbf{\lean{Coord}}\hfill\textsf{\small [load-bearing]}\par\nopagebreak

\begin{leancode}
inductive Coord (S : Shape) (κ : ℕ)
  | const (c : K)
  | known (col : Column κ)
  | committed (c : S.ColumnId) (h : S.τ c.1 = κ)
\end{leancode}

\src{LeanerVM/Protocol/Spine/Instance.lean:86-89 at b435631}

One coordinate of a boundary tuple: a constant, a column both parties know (index, bytecode), or a committed column of the block's height. \emph{Depends on:} \texttt{Shape}, \texttt{Column}.

\par\medskip\noindent\textbf{\lean{BoundaryBlock}}\hfill\textsf{\small [load-bearing]}\par\nopagebreak

\begin{leancode}
structure BoundaryBlock (S : Shape) where
  κ : ℕ
  side : Side
  coords : Vector (Coord S κ) 16
\end{leancode}

\src{LeanerVM/Protocol/Spine/Instance.lean:92-98 at b435631}

\texttt{2\textasciicircum{}κ} boundary tuples on one side of the bus, given coordinate by coordinate (sixteen, separator first). \emph{Depends on:} \texttt{Coord}, \texttt{Side}.

\par\medskip\noindent\textbf{\lean{PublicLine}}\hfill\textsf{\small [load-bearing]}\par\nopagebreak

\begin{leancode}
structure PublicLine (S : Shape) where
  col : S.ColumnId
  cell0 : K
  cell1 : K
  sent : Bool
  pos : 0 < S.τ col.1
\end{leancode}

\src{LeanerVM/Protocol/Spine/Instance.lean:102-114 at b435631}

A column whose cells 0 and 1 the public statement fixes; \texttt{sent} says whether the proof carries its value (transcript only, not the relation); \texttt{pos} says the column has a cell 1. \emph{Depends on:} \texttt{Shape}.

\par\medskip\noindent\textbf{\lean{M3Instance}}\hfill\textsf{\small [load-bearing]}\par\nopagebreak

\begin{leancode}
structure M3Instance extends Shape where
  Stmt : Type
  constraints : (j : Fin ntab) → List (CMvPolynomial (width j) K)
  flushes : (j : Fin ntab) → List (Side × Vector (CMvPolynomial (width j) K) 16)
  d : ℕ
  constraints_degree : ∀ j, ∀ C ∈ constraints j, C.totalDegree ≤ d
  flushes_degree : ∀ j, ∀ f ∈ flushes j, ∀ k, (f.2.get k).totalDegree ≤ d
  counts : (j : Fin ntab) → List (Fin (width j))
  boundary : List (BoundaryBlock toShape)
  μ : ℕ
  layout : Layout μ toShape.ColumnId (fun c ↦ τ c.1)
  publicLines : Stmt → List (PublicLine toShape)
  aux : Column μ → Prop
  decAux : DecidablePred aux
\end{leancode}

\src{LeanerVM/Protocol/Spine/Instance.lean:119-148 at b435631}

Everything the verifier reads about an arithmetization: the statement type, per table its constraints and flushes (polynomials over \texttt{K} of the row), a degree bound \texttt{d} with the two proofs, the count columns, the boundary blocks, the stack height \texttt{μ} and layout, the public lines of a statement, and the auxiliary predicate \texttt{aux} with its decision procedure. Trusted data (section F). \emph{Depends on:} \texttt{Shape}, \texttt{Layout}, \texttt{BoundaryBlock}, \texttt{PublicLine}, CompPoly \texttt{CMvPolynomial}, \texttt{totalDegree}.

\par\medskip\noindent\textbf{\lean{M3Instance.κ}}\hfill\textsf{\small [load-bearing]}\par\nopagebreak

\begin{leancode}
abbrev κ (c : I.ColumnId) : ℕ := I.τ c.1
\end{leancode}

\src{LeanerVM/Protocol/Spine/Instance.lean:157-157 at b435631}

Log-height of a column: its table's. \emph{Depends on:} \texttt{M3Instance}.

\par\medskip\noindent\textbf{\lean{M3Instance.column}}\hfill\textsf{\small [load-bearing]}\par\nopagebreak

\begin{leancode}
def column (q : Column I.μ) (c : I.ColumnId) : Column (I.κ c) := I.layout.read q c
\end{leancode}

\src{LeanerVM/Protocol/Spine/Instance.lean:160-160 at b435631}

One column of the stack, through the layout. \emph{Depends on:} \texttt{Layout.\allowbreak{}read}.

\par\medskip\noindent\textbf{\lean{M3Instance.row}}\hfill\textsf{\small [load-bearing]}\par\nopagebreak

\begin{leancode}
def row (q : Column I.μ) (j : Fin I.ntab) (x : Fin (2 ^ I.τ j)) : Fin (I.width j) → K :=
  fun i ↦ (I.column q ⟨j, i⟩).values.get x
\end{leancode}

\src{LeanerVM/Protocol/Spine/Instance.lean:163-164 at b435631}

Row \texttt{x} of table \texttt{j}: the function giving each column's cell \texttt{x}. \emph{Depends on:} \texttt{column}.

\par\medskip\noindent\textbf{\lean{M3Instance.coordCell}}\hfill\textsf{\small [load-bearing]}\par\nopagebreak

\begin{leancode}
def coordCell (q : Column I.μ) {κ : ℕ} : Coord I.toShape κ → Fin (2 ^ κ) → K
  | .const c, _ => c
  | .known col, x => col.values.get x
  | .committed c h, x => (I.column q c).values.get (Fin.cast (congrArg (2 ^ ·) h.symm) x)
\end{leancode}

\src{LeanerVM/Protocol/Spine/Instance.lean:167-170 at b435631}

The value of a boundary coordinate at row \texttt{x} of its block. \emph{Depends on:} \texttt{Coord}, \texttt{column}.

\par\medskip\noindent\textbf{\lean{M3Instance.flushTuples}}\hfill\textsf{\small [load-bearing]}\par\nopagebreak

\begin{leancode}
def flushTuples (q : Column I.μ) (s : Side) : List (Vector K 16) :=
  (List.finRange I.ntab).flatMap fun j ↦
    ((I.flushes j).filter fun f ↦ decide (f.1 = s)).flatMap fun f ↦
      (List.finRange (2 ^ I.τ j)).map fun x ↦ f.2.map fun P ↦ P.eval (I.row q j x)
\end{leancode}

\src{LeanerVM/Protocol/Spine/Instance.lean:174-177 at b435631}

The tuples of one side flushed by the tables: for every table, every flush of that side, every row, the sixteen coordinate polynomials evaluated on the row. \emph{Depends on:} \texttt{row}, CompPoly \texttt{CMvPolynomial.\allowbreak{}eval}.

\par\medskip\noindent\textbf{\lean{M3Instance.boundaryTuples}}\hfill\textsf{\small [load-bearing]}\par\nopagebreak

\begin{leancode}
def boundaryTuples (q : Column I.μ) (s : Side) : List (Vector K 16) :=
  (I.boundary.filter fun b ↦ decide (b.side = s)).flatMap fun b ↦
    (List.finRange (2 ^ b.κ)).map fun x ↦ b.coords.map fun co ↦ I.coordCell q co x
\end{leancode}

\src{LeanerVM/Protocol/Spine/Instance.lean:180-182 at b435631}

The tuples of one side on the boundary blocks: for every block of that side, every row. \emph{Depends on:} \texttt{coordCell}.

\par\medskip\noindent\textbf{\lean{M3Instance.tuples}}\hfill\textsf{\small [load-bearing]}\par\nopagebreak

\begin{leancode}
def tuples (q : Column I.μ) (s : Side) : List (Vector K 16) :=
  I.flushTuples q s ++ I.boundaryTuples q s
\end{leancode}

\src{LeanerVM/Protocol/Spine/Instance.lean:185-186 at b435631}

Every tuple of one side: the tables' then the boundary's. \emph{Depends on:} \texttt{flushTuples}, \texttt{boundaryTuples}.

\par\medskip\noindent\textbf{\lean{M3Instance.ConstraintsVanish}}\hfill\textsf{\small [load-bearing]}\par\nopagebreak

\begin{leancode}
def ConstraintsVanish (q : Column I.μ) : Prop :=
  ∀ j, ∀ C ∈ I.constraints j, ∀ x, C.eval (I.row q j x) = 0
\end{leancode}

\src{LeanerVM/Protocol/Spine/Instance.lean:197-198 at b435631}

Every constraint of every table is zero on every row (the zerocheck target). \emph{Depends on:} \texttt{row}.

\par\medskip\noindent\textbf{\lean{M3Instance.Balanced}}\hfill\textsf{\small [load-bearing]}\par\nopagebreak

\begin{leancode}
def Balanced (q : Column I.μ) : Prop := (I.tuples q .push).Perm (I.tuples q .pull)
\end{leancode}

\src{LeanerVM/Protocol/Spine/Instance.lean:201-201 at b435631}

The pushed tuples are a permutation of the pulled tuples: the same multiset, counted in \texttt{ℕ}. \emph{Depends on:} \texttt{tuples}, Mathlib \texttt{List.\allowbreak{}Perm}.

\par\medskip\noindent\textbf{\lean{M3Instance.CountsNonzero}}\hfill\textsf{\small [load-bearing]}\par\nopagebreak

\begin{leancode}
def CountsNonzero (q : Column I.μ) : Prop :=
  ∀ j, ∀ i ∈ I.counts j, ∀ x, (I.column q ⟨j, i⟩).values.get x ≠ 0
\end{leancode}

\src{LeanerVM/Protocol/Spine/Instance.lean:204-205 at b435631}

Every cell of every count column is nonzero (the count product). \emph{Depends on:} \texttt{column}.

\par\medskip\noindent\textbf{\lean{M3Instance.PublicLinesHold}}\hfill\textsf{\small [load-bearing]}\par\nopagebreak

\begin{leancode}
def PublicLinesHold (input : I.Stmt) (q : Column I.μ) : Prop :=
  ∀ l ∈ I.publicLines input,
    (I.column q l.col).values.get ⟨0, Nat.two_pow_pos _⟩ = l.cell0 ∧
      (I.column q l.col).values.get ⟨1, Nat.one_lt_two_pow l.pos.ne'⟩ = l.cell1
\end{leancode}

\src{LeanerVM/Protocol/Spine/Instance.lean:208-211 at b435631}

Cells 0 and 1 of every line the statement names hold the line's values. \emph{Depends on:} \texttt{column}, \texttt{PublicLine}.

\par\medskip\noindent\textbf{\lean{M3Holds}}\hfill\textsf{\small [load-bearing]}\par\nopagebreak

\begin{leancode}
def M3Holds (I : M3Instance) (input : I.Stmt) (q : Column I.μ) : Prop :=
  I.ConstraintsVanish q ∧ I.Balanced q ∧ I.CountsNonzero q ∧ I.PublicLinesHold input q ∧ I.aux q
\end{leancode}

\src{LeanerVM/Protocol/Spine/Instance.lean:219-220 at b435631}

The relation: the five clauses. What the verifier establishes about the committed stack. \emph{Depends on:} the five clauses.

\par\medskip\noindent\textbf{\lean{TheOracle}}\hfill\textsf{\small [load-bearing]}\par\nopagebreak

\begin{leancode}
abbrev TheOracle (I : M3Instance) : Fin 1 → Type := OneOracle (Column I.μ)
\end{leancode}

\src{LeanerVM/Protocol/Spine/Instance.lean:226-226 at b435631}

The one oracle from the commit on: the stack. \emph{Depends on:} \texttt{OneOracle}, \texttt{Column}.

\par\medskip\noindent\textbf{\lean{M3Rel}}\hfill\textsf{\small [load-bearing]}\par\nopagebreak

\begin{leancode}
def M3Rel (I : M3Instance) : Set ((I.Stmt × (∀ i, NoOracle i)) × Column I.μ) :=
  {p | M3Holds I p.1.1 p.2}
\end{leancode}

\src{LeanerVM/Protocol/Spine/Instance.lean:230-231 at b435631}

\texttt{M3Holds} as ArkLib states relations: statement = public input with no oracle, witness = the stack. \emph{Depends on:} \texttt{M3Holds}, \texttt{NoOracle}.

\subsection{\texttt{LeanerVM/Protocol/Spine/Seams.lean}}\label{sec:gen-cat-spine-Seams}

\par\medskip\noindent\textbf{\lean{ColumnClaim}}\hfill\textsf{\small [load-bearing]}\par\nopagebreak

\begin{leancode}
structure ColumnClaim (I : M3Instance) where
  col : I.ColumnId
  point : Vector E (I.κ col)
  value : E
\end{leancode}

\src{LeanerVM/Protocol/Spine/Seams.lean:57-63 at b435631}

One column's extension at a point of \texttt{E\textasciicircum{}κ} equals a value. \emph{Depends on:} \texttt{M3Instance}.

\par\medskip\noindent\textbf{\lean{ColumnClaim.Holds}}\hfill\textsf{\small [load-bearing]}\par\nopagebreak

\begin{leancode}
def ColumnClaim.Holds {I : M3Instance} (q : Column I.μ) (c : ColumnClaim I) : Prop :=
  CMlPolynomialEval.eval₂Mle (I.column q c.col).values (algebraMap K E) c.point = c.value
\end{leancode}

\src{LeanerVM/Protocol/Spine/Seams.lean:66-67 at b435631}

The claim is true of the stack, reading the column through the layout. \emph{Depends on:} \texttt{column}, \texttt{eval₂Mle}.

\par\medskip\noindent\textbf{\lean{VirtualTerm}}\hfill\textsf{\small [load-bearing]}\par\nopagebreak

\begin{leancode}
structure VirtualTerm (I : M3Instance) where
  weight : E
  j : Fin I.ntab
  poly : CMvPolynomial (I.width j) K
  point : Vector E (I.τ j)
\end{leancode}

\src{LeanerVM/Protocol/Spine/Seams.lean:71-79 at b435631}

One term of a linear claim: a weight in \texttt{E}, a table, a polynomial of its row over \texttt{K}, and the point the virtual table (the polynomial on every row) is extended to. \emph{Depends on:} \texttt{M3Instance}.

\par\medskip\noindent\textbf{\lean{VirtualTerm.table}}\hfill\textsf{\small [load-bearing]}\par\nopagebreak

\begin{leancode}
def VirtualTerm.table {I : M3Instance} (q : Column I.μ) (t : VirtualTerm I) :
    CMlPolynomialEval E (I.τ t.j) :=
  Vector.ofFn fun x ↦ ofK (t.poly.eval (I.row q t.j x))
\end{leancode}

\src{LeanerVM/Protocol/Spine/Seams.lean:82-84 at b435631}

The virtual table of a term: its polynomial evaluated on every row of its table, lifted to \texttt{E}. \emph{Depends on:} \texttt{row}.

\par\medskip\noindent\textbf{\lean{VirtualTerm.eval}}\hfill\textsf{\small [load-bearing]}\par\nopagebreak

\begin{leancode}
def VirtualTerm.eval {I : M3Instance} (q : Column I.μ) (t : VirtualTerm I) : E :=
  t.weight * CMlPolynomialEval.evalMle (t.table q) t.point
\end{leancode}

\src{LeanerVM/Protocol/Spine/Seams.lean:87-88 at b435631}

The term's value: weight times the virtual table's extension at the point. \emph{Depends on:} \texttt{table}, CompPoly \texttt{evalMle}.

\par\medskip\noindent\textbf{\lean{LinearClaim}}\hfill\textsf{\small [load-bearing]}\par\nopagebreak

\begin{leancode}
structure LinearClaim (I : M3Instance) where
  terms : List (VirtualTerm I)
  value : E
\end{leancode}

\src{LeanerVM/Protocol/Spine/Seams.lean:91-95 at b435631}

A weighted sum of virtual-table extensions equals a value (a zerocheck claim is one term of weight 1 and value 0; a bus form is a list of terms). \emph{Depends on:} \texttt{VirtualTerm}.

\par\medskip\noindent\textbf{\lean{LinearClaim.Holds}}\hfill\textsf{\small [load-bearing]}\par\nopagebreak

\begin{leancode}
def LinearClaim.Holds {I : M3Instance} (q : Column I.μ) (c : LinearClaim I) : Prop :=
  (c.terms.map fun t ↦ t.eval q).sum = c.value
\end{leancode}

\src{LeanerVM/Protocol/Spine/Seams.lean:98-99 at b435631}

The sum of the terms' values is the claimed value. \emph{Depends on:} \texttt{VirtualTerm.\allowbreak{}eval}.

\par\medskip\noindent\textbf{\lean{Weight}}\hfill\textsf{\small [load-bearing]}\par\nopagebreak

\begin{leancode}
structure Weight (μ : ℕ) where
  onCube : CMlPolynomialEval E μ
  mle : Vector E μ → E
  mle_eq : ∀ r, mle r = CMlPolynomialEval.evalMle onCube r
\end{leancode}

\src{LeanerVM/Protocol/Spine/Seams.lean:103-109 at b435631}

A weight on the stack (specification Definition 3.13): its values on the cube, an evaluator for its extension, and their agreement. \emph{Depends on:} CompPoly \texttt{evalMle}.

\par\medskip\noindent\textbf{\lean{Weight.pair}}\hfill\textsf{\small [load-bearing]}\par\nopagebreak

\begin{leancode}
def Weight.pair {μ : ℕ} (W : Weight μ) (q : Column μ) : E :=
  ∑ i : Fin (2 ^ μ), W.onCube.get i * ofK (q.values.get i)
\end{leancode}

\src{LeanerVM/Protocol/Spine/Seams.lean:112-113 at b435631}

\texttt{Σ\_\allowbreak{}x W(\allowbreak{}x)·q(\allowbreak{}x)} over the cube. \emph{Depends on:} \texttt{Weight}, \texttt{Column}.

\par\medskip\noindent\textbf{\lean{WeightedClaim}}\hfill\textsf{\small [load-bearing]}\par\nopagebreak

\begin{leancode}
structure WeightedClaim (I : M3Instance) where
  weight : Weight I.μ
  value : E
\end{leancode}

\src{LeanerVM/Protocol/Spine/Seams.lean:116-120 at b435631}

The pairing of a weight with the stack equals a value. \emph{Depends on:} \texttt{Weight}.

\par\medskip\noindent\textbf{\lean{WeightedClaim.Holds}}\hfill\textsf{\small [load-bearing]}\par\nopagebreak

\begin{leancode}
def WeightedClaim.Holds {I : M3Instance} (q : Column I.μ) (c : WeightedClaim I) : Prop :=
  c.weight.pair q = c.value
\end{leancode}

\src{LeanerVM/Protocol/Spine/Seams.lean:123-124 at b435631}

True of the stack when the pairing is the value. \emph{Depends on:} \texttt{Weight.\allowbreak{}pair}.

\par\medskip\noindent\textbf{\lean{BusOut}}\hfill\textsf{\small [load-bearing]}\par\nopagebreak

\begin{leancode}
structure BusOut (I : M3Instance) where
  linear : List (LinearClaim I)
  columns : List (ColumnClaim I)
\end{leancode}

\src{LeanerVM/Protocol/Spine/Seams.lean:130-134 at b435631}

What the bus phase hands on: linear claims and column claims. \emph{Depends on:} \texttt{LinearClaim}, \texttt{ColumnClaim}.

\par\medskip\noindent\textbf{\lean{TableOut}}\hfill\textsf{\small [load-bearing]}\par\nopagebreak

\begin{leancode}
structure TableOut (I : M3Instance) where
  columns : List (ColumnClaim I)
\end{leancode}

\src{LeanerVM/Protocol/Spine/Seams.lean:137-139 at b435631}

What the table sumcheck hands on: column claims. \emph{Depends on:} \texttt{ColumnClaim}.

\par\medskip\noindent\textbf{\lean{PubOut}}\hfill\textsf{\small [load-bearing]}\par\nopagebreak

\begin{leancode}
structure PubOut (I : M3Instance) where
  columns : List (ColumnClaim I)
\end{leancode}

\src{LeanerVM/Protocol/Spine/Seams.lean:142-144 at b435631}

What the public-input phase hands on: column claims (the same type as \texttt{TableOut}). \emph{Depends on:} \texttt{ColumnClaim}.

\par\medskip\noindent\textbf{\lean{FlockOut}}\hfill\textsf{\small [load-bearing]}\par\nopagebreak

\begin{leancode}
structure FlockOut (I : M3Instance) where
  columns : List (ColumnClaim I)
  weighted : List (WeightedClaim I)
\end{leancode}

\src{LeanerVM/Protocol/Spine/Seams.lean:147-151 at b435631}

What the Flock phase hands on: column claims and weighted claims. \emph{Depends on:} \texttt{ColumnClaim}, \texttt{WeightedClaim}.

\par\medskip\noindent\textbf{\lean{theStack}}\hfill\textsf{\small [load-bearing]}\par\nopagebreak

\begin{leancode}
abbrev theStack {I : M3Instance} (o : ∀ i, TheOracle I i) : Column I.μ := o 0
\end{leancode}

\src{LeanerVM/Protocol/Spine/Seams.lean:159-159 at b435631}

The stack behind the one oracle. \emph{Depends on:} \texttt{TheOracle}.

\par\medskip\noindent\textbf{\lean{Seam.of}}\hfill\textsf{\small [load-bearing]}\par\nopagebreak

\begin{leancode}
def of {S : Type} (P : S → Column I.μ → Prop) : Set ((S × ∀ i, TheOracle I i) × Unit) :=
  {p | P p.1.1 (theStack p.1.2)}
\end{leancode}

\src{LeanerVM/Protocol/Spine/Seams.lean:166-167 at b435631}

A seam from a predicate on the statement and the stack: the set of \texttt{(\allowbreak{}(\allowbreak{}statement,\allowbreak{} oracles),\allowbreak{} (\allowbreak{}))} where it holds of the one oracle. \emph{Depends on:} \texttt{theStack}.

\par\medskip\noindent\textbf{\lean{Seam.commit}}\hfill\textsf{\small [load-bearing]}\par\nopagebreak

\begin{leancode}
def commit := of I (M3Holds I)
\end{leancode}

\src{LeanerVM/Protocol/Spine/Seams.lean:171-171 at b435631}

After the commit: \texttt{M3Holds} of the oracle itself. \emph{Depends on:} \texttt{of}, \texttt{M3Holds}.

\par\medskip\noindent\textbf{\lean{Seam.bus}}\hfill\textsf{\small [load-bearing]}\par\nopagebreak

\begin{leancode}
def bus := of I fun (s : I.Stmt × BusOut I) q ↦ (∀ c ∈ s.2.linear, c.Holds q) ∧
  (∀ c ∈ s.2.linear, ∀ t ∈ c.terms, t.poly.totalDegree ≤ I.d) ∧
  (∀ c ∈ s.2.columns, c.Holds q) ∧ I.PublicLinesHold s.1 q ∧ I.aux q
\end{leancode}

\src{LeanerVM/Protocol/Spine/Seams.lean:175-177 at b435631}

After the bus phase: every linear claim holds, every term is within the degree bound, every column claim holds, the lines hold, \texttt{aux} holds. \emph{Depends on:} \texttt{of}, the claims, \texttt{PublicLinesHold}.

\par\medskip\noindent\textbf{\lean{Seam.table}}\hfill\textsf{\small [load-bearing]}\par\nopagebreak

\begin{leancode}
def table := of I fun (s : I.Stmt × TableOut I) q ↦
  (∀ c ∈ s.2.columns, c.Holds q) ∧ I.PublicLinesHold s.1 q ∧ I.aux q
\end{leancode}

\src{LeanerVM/Protocol/Spine/Seams.lean:180-181 at b435631}

After the table sumcheck: column claims, lines, \texttt{aux}. \emph{Depends on:} \texttt{of}.

\par\medskip\noindent\textbf{\lean{Seam.pub}}\hfill\textsf{\small [load-bearing]}\par\nopagebreak

\begin{leancode}
def pub := of I fun (s : I.Stmt × PubOut I) q ↦ (∀ c ∈ s.2.columns, c.Holds q) ∧ I.aux q
\end{leancode}

\src{LeanerVM/Protocol/Spine/Seams.lean:184-184 at b435631}

After the public-input phase: column claims and \texttt{aux}. \emph{Depends on:} \texttt{of}.

\par\medskip\noindent\textbf{\lean{Seam.flock}}\hfill\textsf{\small [load-bearing]}\par\nopagebreak

\begin{leancode}
def flock := of I fun (s : I.Stmt × FlockOut I) q ↦
  (∀ c ∈ s.2.columns, c.Holds q) ∧ ∀ c ∈ s.2.weighted, c.Holds q
\end{leancode}

\src{LeanerVM/Protocol/Spine/Seams.lean:187-188 at b435631}

After the Flock phase: column claims and weighted claims. \emph{Depends on:} \texttt{of}.

\par\medskip\noindent\textbf{\lean{Seam.done}}\hfill\textsf{\small [load-bearing]}\par\nopagebreak

\begin{leancode}
def done := of I fun (_ : Unit) _ ↦ True
\end{leancode}

\src{LeanerVM/Protocol/Spine/Seams.lean:191-191 at b435631}

After the opening: nothing (the whole set). \emph{Depends on:} \texttt{of}.

\subsection{\texttt{LeanerVM/Protocol/Spine/Phase.lean}}\label{sec:gen-cat-spine-Phase}

\par\medskip\noindent\textbf{\lean{Phase.Def}}\hfill\textsf{\small [interface]}\par\nopagebreak

\begin{leancode}
abbrev Def (StmtIn StmtOut : Type) : Type 1 :=
  Component.Def StmtIn (TheOracle I) Unit StmtOut (TheOracle I) Unit
\end{leancode}

\src{LeanerVM/Protocol/Spine/Phase.lean:37-38 at b435631}

A phase: a component from a statement to a statement over the stack, with trivial witnesses. \emph{Depends on:} \texttt{Component.\allowbreak{}Def}, \texttt{TheOracle}.

\par\medskip\noindent\textbf{\lean{Phase.Complete}}\hfill\textsf{\small [interface]}\par\nopagebreak

\begin{leancode}
abbrev Complete {StmtIn StmtOut : Type} (D : Def I StmtIn StmtOut)
    (relIn : Set ((StmtIn × ∀ i, TheOracle I i) × Unit))
    (relOut : Set ((StmtOut × ∀ i, TheOracle I i) × Unit)) : Type :=
  Component.Complete D relIn relOut
\end{leancode}

\src{LeanerVM/Protocol/Spine/Phase.lean:41-44 at b435631}

Its completeness half against two seams. \emph{Depends on:} \texttt{Component.\allowbreak{}Complete}.

\par\medskip\noindent\textbf{\lean{Phase.Security}}\hfill\textsf{\small [interface]}\par\nopagebreak

\begin{leancode}
abbrev Security {StmtIn StmtOut : Type} (D : Def I StmtIn StmtOut)
    (relIn : Set ((StmtIn × ∀ i, TheOracle I i) × Unit))
    (relOut : Set ((StmtOut × ∀ i, TheOracle I i) × Unit)) : Type 1 :=
  Component.Security D relIn relOut
\end{leancode}

\src{LeanerVM/Protocol/Spine/Phase.lean:47-50 at b435631}

Its security half against two seams. \emph{Depends on:} \texttt{Component.\allowbreak{}Security}.

\par\medskip\noindent\textbf{\lean{Phase.passThrough}}\hfill\textsf{\small [interface]}\par\nopagebreak

\begin{leancode}
abbrev passThrough (f : StmtIn → StmtOut) : Def I StmtIn StmtOut :=
  Component.passThrough (TheOracle I) f
\end{leancode}

\src{LeanerVM/Protocol/Spine/Phase.lean:55-56 at b435631}

The pass-through phase. \emph{Depends on:} \texttt{Component.\allowbreak{}passThrough}.

\par\medskip\noindent\textbf{\lean{Phase.passThroughComplete}}\hfill\textsf{\small [interface]}\par\nopagebreak

\begin{leancode}
def passThroughComplete (f : StmtIn → StmtOut)
    {relIn : Set ((StmtIn × ∀ i, TheOracle I i) × Unit)}
    {relOut : Set ((StmtOut × ∀ i, TheOracle I i) × Unit)}
    (h : ∀ s o, ((s, o), ()) ∈ relIn → ((f s, o), ()) ∈ relOut) :
    Complete I (passThrough I f) relIn relOut :=
  Component.passThroughComplete (TheOracle I) f fun s o _ hin ↦ h s o hin
\end{leancode}

\src{LeanerVM/Protocol/Spine/Phase.lean:59-64 at b435631}

Its completeness when the map carries one seam into the other. \emph{Depends on:} \texttt{C\allowbreak{}o\allowbreak{}m\allowbreak{}p\allowbreak{}o\allowbreak{}n\allowbreak{}e\allowbreak{}n\allowbreak{}t\allowbreak{}.\allowbreak{}p\allowbreak{}a\allowbreak{}s\allowbreak{}s\allowbreak{}T\allowbreak{}h\allowbreak{}r\allowbreak{}o\allowbreak{}u\allowbreak{}g\allowbreak{}h\allowbreak{}C\allowbreak{}o\allowbreak{}m\allowbreak{}p\allowbreak{}l\allowbreak{}e\allowbreak{}t\allowbreak{}e}.

\par\medskip\noindent\textbf{\lean{Phase.passThroughSecurity}}\hfill\textsf{\small [interface]}\par\nopagebreak

\begin{leancode}
def passThroughSecurity (f : StmtIn → StmtOut)
    {relIn : Set ((StmtIn × ∀ i, TheOracle I i) × Unit)}
    {relOut : Set ((StmtOut × ∀ i, TheOracle I i) × Unit)}
    (hc : ∀ s o, ((s, o), ()) ∈ relIn → ((f s, o), ()) ∈ relOut)
    (h : ∀ s o, ((f s, o), ()) ∈ relOut → ((s, o), ()) ∈ relIn) :
    Security I (passThrough I f) relIn relOut :=
  Component.passThroughSecurity (TheOracle I) f (fun s o _ ↦ h s o) fun s o _ ↦ hc s o
\end{leancode}

\src{LeanerVM/Protocol/Spine/Phase.lean:68-74 at b435631}

Its security when the map also reflects the second seam into the first (D.3 (a): impossible at the bus seam). \emph{Depends on:} \texttt{C\allowbreak{}o\allowbreak{}m\allowbreak{}p\allowbreak{}o\allowbreak{}n\allowbreak{}e\allowbreak{}n\allowbreak{}t\allowbreak{}.\allowbreak{}p\allowbreak{}a\allowbreak{}s\allowbreak{}s\allowbreak{}T\allowbreak{}h\allowbreak{}r\allowbreak{}o\allowbreak{}u\allowbreak{}g\allowbreak{}h\allowbreak{}S\allowbreak{}e\allowbreak{}c\allowbreak{}u\allowbreak{}r\allowbreak{}i\allowbreak{}t\allowbreak{}y}.

\subsection{\texttt{LeanerVM/Protocol/Spine/Compose.lean}}\label{sec:gen-cat-spine-Compose}

\par\medskip\noindent\textbf{\lean{commitSpec}}\hfill\textsf{\small [load-bearing]}\par\nopagebreak

\begin{leancode}
abbrev commitSpec : ProtocolSpec 1 := Component.sendSpec (Column I.μ)
\end{leancode}

\src{LeanerVM/Protocol/Spine/Compose.lean:51-51 at b435631}

The commit phase's schedule: one prover message, the stack. \emph{Depends on:} \texttt{sendSpec}.

\par\medskip\noindent\textbf{\lean{commitDef}}\hfill\textsf{\small [load-bearing]}\par\nopagebreak

\begin{leancode}
abbrev commitDef : Component.Def I.Stmt NoOracle (Column I.μ) I.Stmt (TheOracle I) Unit :=
  Component.sendOracle I.Stmt (Column I.μ)
\end{leancode}

\src{LeanerVM/Protocol/Spine/Compose.lean:55-56 at b435631}

The commit phase: the send-oracle component at the stack. \emph{Depends on:} \texttt{sendOracle}.

\par\medskip\noindent\textbf{\lean{commitExtractor}}\hfill\textsf{\small [load-bearing (named form)]}\par\nopagebreak

\begin{leancode}
abbrev commitExtractor := Component.sendExtractor (S := I.Stmt) (M := Column I.μ)
\end{leancode}

\src{LeanerVM/Protocol/Spine/Compose.lean:59-59 at b435631}

Its extractor: read the stack off the message. \emph{Depends on:} \texttt{sendExtractor}.

\par\medskip\noindent\textbf{\lean{commitComplete}}\hfill\textsf{\small [load-bearing (named form)]}\par\nopagebreak

\begin{leancode}
def commitComplete : Component.Complete (commitDef I) (M3Rel I) (Seam.commit I) :=
  Component.sendOracleComplete (M3Rel I)
\end{leancode}

\src{LeanerVM/Protocol/Spine/Compose.lean:63-64 at b435631}

Its completeness from \texttt{M3Rel} to \texttt{Seam.\allowbreak{}commit}. \emph{Depends on:} \texttt{sendOracleComplete}.

\par\medskip\noindent\textbf{\lean{commitSecurity}}\hfill\textsf{\small [load-bearing (named form)]}\par\nopagebreak

\begin{leancode}
def commitSecurity : Component.Security (commitDef I) (M3Rel I) (Seam.commit I) :=
  Component.sendOracleSecurity (M3Rel I)
\end{leancode}

\src{LeanerVM/Protocol/Spine/Compose.lean:67-68 at b435631}

Its security at error zero. \emph{Depends on:} \texttt{sendOracleSecurity}.

\par\medskip\noindent\textbf{\lean{Phases}}\hfill\textsf{\small [load-bearing]}\par\nopagebreak

\begin{leancode}
structure Phases where
  bus : Phase.Def I I.Stmt (I.Stmt × BusOut I)
  table : Phase.Def I (I.Stmt × BusOut I) (I.Stmt × TableOut I)
  pub : Phase.Def I (I.Stmt × TableOut I) (I.Stmt × PubOut I)
  flock : Phase.Def I (I.Stmt × PubOut I) (I.Stmt × FlockOut I)
  opening : Phase.Def I (I.Stmt × FlockOut I) Unit
\end{leancode}

\src{LeanerVM/Protocol/Spine/Compose.lean:73-84 at b435631}

The five phases after the commit, by their statement types. \emph{Depends on:} \texttt{Phase.\allowbreak{}Def}, the four outputs.

\par\medskip\noindent\textbf{\lean{Phases.toDef}}\hfill\textsf{\small [load-bearing]}\par\nopagebreak

\begin{leancode}
def Phases.toDef (P : Phases I) :
    Component.Def I.Stmt NoOracle (Column I.μ) Unit (TheOracle I) Unit :=
  (((((commitDef I).append P.bus).append P.table).append P.pub).append P.flock).append P.opening
\end{leancode}

\src{LeanerVM/Protocol/Spine/Compose.lean:89-91 at b435631}

The whole protocol as one component: the commit, then five appends. \emph{Depends on:} \texttt{commitDef}, \texttt{Def.\allowbreak{}append}.

\par\medskip\noindent\textbf{\lean{Phases.Complete}}\hfill\textsf{\small [load-bearing]}\par\nopagebreak

\begin{leancode}
structure Phases.Complete (P : Phases I) where
  bus : Phase.Complete I P.bus (Seam.commit I) (Seam.bus I)
  table : Phase.Complete I P.table (Seam.bus I) (Seam.table I)
  pub : Phase.Complete I P.pub (Seam.table I) (Seam.pub I)
  flock : Phase.Complete I P.flock (Seam.pub I) (Seam.flock I)
  opening : Phase.Complete I P.opening (Seam.flock I) (Seam.done I)
\end{leancode}

\src{LeanerVM/Protocol/Spine/Compose.lean:94-104 at b435631}

The five completeness halves against the seams. \emph{Depends on:} \texttt{Phase.\allowbreak{}Complete}, the seams.

\par\medskip\noindent\textbf{\lean{Phases.Security}}\hfill\textsf{\small [load-bearing]}\par\nopagebreak

\begin{leancode}
structure Phases.Security (P : Phases I) where
  bus : Phase.Security I P.bus (Seam.commit I) (Seam.bus I)
  table : Phase.Security I P.table (Seam.bus I) (Seam.table I)
  pub : Phase.Security I P.pub (Seam.table I) (Seam.pub I)
  flock : Phase.Security I P.flock (Seam.pub I) (Seam.flock I)
  opening : Phase.Security I P.opening (Seam.flock I) (Seam.done I)
\end{leancode}

\src{LeanerVM/Protocol/Spine/Compose.lean:113-123 at b435631}

The five security halves against the seams. \emph{Depends on:} \texttt{Phase.\allowbreak{}Security}, the seams.

\par\medskip\noindent\textbf{\lean{Phases.Security.toDef}}\hfill\textsf{\small [load-bearing (named form)]}\par\nopagebreak

\begin{leancode}
def Phases.Security.toDef {P : Phases I} (S : P.Security) :
    Component.Security P.toDef (M3Rel I) (Seam.done I) :=
  (((((commitSecurity I).append S.bus).append S.table).append S.pub).append S.flock).append
    S.opening
\end{leancode}

\src{LeanerVM/Protocol/Spine/Compose.lean:126-129 at b435631}

The composed security: the commit's, then five appends. \emph{Depends on:} \texttt{commitSecurity}, \texttt{Security.\allowbreak{}append}.

\par\medskip\noindent\textbf{\lean{leanVmPiop}}\hfill\textsf{\small [load-bearing]}\par\nopagebreak

\begin{leancode}
def leanVmPiop (P : Phases I) :
    OracleReduction []ₒ I.Stmt NoOracle (Column I.μ) Unit (TheOracle I) Unit P.toDef.pSpec :=
  P.toDef.red
\end{leancode}

\src{LeanerVM/Protocol/Spine/Compose.lean:134-136 at b435631}

The oracle protocol: the composed reduction. \emph{Depends on:} \texttt{Phases.\allowbreak{}toDef}.

\par\medskip\noindent\textbf{\lean{leanVmVerifier}}\hfill\textsf{\small [load-bearing]}\par\nopagebreak

\begin{leancode}
def leanVmVerifier (P : Phases I) :
    OracleVerifier []ₒ I.Stmt NoOracle Unit (TheOracle I) P.toDef.pSpec :=
  P.toDef.red.verifier
\end{leancode}

\src{LeanerVM/Protocol/Spine/Compose.lean:139-141 at b435631}

Its verifier. \emph{Depends on:} \texttt{Phases.\allowbreak{}toDef}.

\par\medskip\noindent\textbf{\lean{leanVmProver}}\hfill\textsf{\small [interface]}\par\nopagebreak

\begin{leancode}
def leanVmProver (P : Phases I) :
    OracleProver []ₒ I.Stmt NoOracle (Column I.μ) Unit (TheOracle I) Unit P.toDef.pSpec :=
  P.toDef.red.prover
\end{leancode}

\src{LeanerVM/Protocol/Spine/Compose.lean:144-146 at b435631}

Its honest prover (not in either theorem's statement; \texttt{leanVmPiop} carries it). \emph{Depends on:} \texttt{Phases.\allowbreak{}toDef}.

\par\medskip\noindent\textbf{\lean{piopError}}\hfill\textsf{\small [load-bearing]}\par\nopagebreak

\begin{leancode}
def piopError (P : Phases I) : P.toDef.pSpec.ChallengeIdx → ℝ≥0 := P.toDef.err
\end{leancode}

\src{LeanerVM/Protocol/Spine/Compose.lean:149-149 at b435631}

Its error per challenge: the phases' declarations side by side (section F). \emph{Depends on:} \texttt{Def.\allowbreak{}append}.

\par\medskip\noindent\textbf{\lean{piopExtractor}}\hfill\textsf{\small [load-bearing (named form)]}\par\nopagebreak

\begin{leancode}
def piopExtractor (P : Phases I) (S : P.Security) :
    Extractor.RoundByRound (OracleSpec.emptySpec.{0, 0}) (I.Stmt × ∀ i, NoOracle i)
      (Column I.μ) Unit P.toDef.pSpec S.toDef.witMid :=
  S.toDef.extractor
\end{leancode}

\src{LeanerVM/Protocol/Spine/Compose.lean:153-156 at b435631}

Its extractor: the composed one (D.5). \emph{Depends on:} \texttt{Phases.\allowbreak{}Security.\allowbreak{}toDef}.

\par\medskip\noindent\textbf{\lean{piop_perfectCompleteness}}\hfill\textsf{\small [master theorem]}\par\nopagebreak

\begin{leancode}
theorem piop_perfectCompleteness (P : Phases I) (C : P.Complete) {σ : Type}
    (init : ProbComp σ) (impl : QueryImpl []ₒ (StateT σ ProbComp)) :
    (leanVmPiop P).perfectCompleteness init impl (M3Rel I) (Seam.done I) :=
\end{leancode}

\src{LeanerVM/Protocol/Spine/Compose.lean:162-165 at b435631}

Given every phase's completeness, the honest prover convinces the verifier with probability one on every \texttt{(\allowbreak{}input,\allowbreak{} q)} with \texttt{M3Holds I input q}. \emph{Depends on:} \texttt{leanVmPiop}, \texttt{M3Rel}, \texttt{Seam.\allowbreak{}done}, ArkLib \texttt{perfectCompleteness}.

\par\medskip\noindent\textbf{\lean{piop_rbrKnowledgeSoundness}}\hfill\textsf{\small [master theorem]}\par\nopagebreak

\begin{leancode}
theorem piop_rbrKnowledgeSoundness (P : Phases I) (S : P.Security) {σ : Type}
    (init : ProbComp σ) (impl : QueryImpl []ₒ (StateT σ ProbComp)) :
    (leanVmVerifier P).toVerifier.rbrKnowledgeSoundnessWorstCaseWith init impl (M3Rel I)
      (Seam.done I) S.toDef.witMid (piopExtractor P S) (S.toDef.kSF init impl)
      (piopError P) :=
\end{leancode}

\src{LeanerVM/Protocol/Spine/Compose.lean:172-177 at b435631}

Given every phase's security, the verifier is round-by-round knowledge sound at the composed error, for the composed extractor and state function. \emph{Depends on:} \texttt{leanVmVerifier}, \texttt{piopExtractor}, \texttt{piopError}, ArkLib \texttt{r\allowbreak{}b\allowbreak{}r\allowbreak{}K\allowbreak{}n\allowbreak{}o\allowbreak{}w\allowbreak{}l\allowbreak{}e\allowbreak{}d\allowbreak{}g\allowbreak{}e\allowbreak{}S\allowbreak{}o\allowbreak{}u\allowbreak{}n\allowbreak{}d\allowbreak{}n\allowbreak{}e\allowbreak{}s\allowbreak{}s\allowbreak{}W\allowbreak{}o\allowbreak{}r\allowbreak{}s\allowbreak{}t\allowbreak{}C\allowbreak{}a\allowbreak{}s\allowbreak{}e\allowbreak{}W\allowbreak{}i\allowbreak{}t\allowbreak{}h}.

\par\medskip\noindent\textbf{\lean{piop_rbrKnowledgeSoundness_exists}}\hfill\textsf{\small [master theorem (existential form)]}\par\nopagebreak

\begin{leancode}
theorem piop_rbrKnowledgeSoundness_exists (P : Phases I) (S : P.Security) {σ : Type}
    (init : ProbComp σ) (impl : QueryImpl []ₒ (StateT σ ProbComp)) :
    (leanVmVerifier P).toVerifier.rbrKnowledgeSoundnessWorstCase init impl (M3Rel I)
      (Seam.done I) (piopError P) :=
\end{leancode}

\src{LeanerVM/Protocol/Spine/Compose.lean:181-185 at b435631}

The same with the extractor and state function forgotten. \emph{Depends on:} \texttt{piop\_\allowbreak{}rbrKnowledgeSoundness}.

\par\medskip\noindent\textbf{Helpers and test fixtures of \file{Compose.lean}.}

{\footnotesize
\begin{longtable}{@{}>{\raggedright\arraybackslash}p{0.300\linewidth}>{\raggedright\arraybackslash}p{0.100\linewidth}>{\raggedright\arraybackslash}p{0.160\linewidth}>{\raggedright\arraybackslash}p{0.360\linewidth}@{}}
\toprule
Declaration & Lines & Class & Meaning \\
\midrule
\endfirsthead
\toprule
Declaration & Lines & Class & Meaning \\
\midrule
\endhead
\bottomrule
\endfoot
\texttt{P\allowbreak{}h\allowbreak{}a\allowbreak{}s\allowbreak{}e\allowbreak{}s\allowbreak{}.\allowbreak{}C\allowbreak{}o\allowbreak{}m\allowbreak{}p\allowbreak{}l\allowbreak{}e\allowbreak{}t\allowbreak{}e\allowbreak{}.\allowbreak{}t\allowbreak{}o\allowbreak{}D\allowbreak{}e\allowbreak{}f} & 107-110 & helper & The composed completeness. \\
\end{longtable}}

\subsection{\texttt{LeanerVM/Protocol/Spine/Toy.lean}}\label{sec:gen-cat-spine-Toy}

\par\medskip\noindent\textbf{Helpers and test fixtures of \file{Toy.lean}.}

{\footnotesize
\begin{longtable}{@{}>{\raggedright\arraybackslash}p{0.300\linewidth}>{\raggedright\arraybackslash}p{0.100\linewidth}>{\raggedright\arraybackslash}p{0.160\linewidth}>{\raggedright\arraybackslash}p{0.360\linewidth}@{}}
\toprule
Declaration & Lines & Class & Meaning \\
\midrule
\endfirsthead
\toprule
Declaration & Lines & Class & Meaning \\
\midrule
\endhead
\bottomrule
\endfoot
\texttt{shape} & 36-36 & test fixture & One table of height two and width three. \\
\texttt{slice} & 42-44 & test fixture & Column \texttt{i} is the slice at cells \texttt{2i,\allowbreak{} 2i+1}. \\
\texttt{extend} & 47-48 & test fixture & The selector of column \texttt{i}: \texttt{(\allowbreak{}z\allowbreak{},\allowbreak{} i\allowbreak{} m\allowbreak{}o\allowbreak{}d\allowbreak{} 2\allowbreak{},\allowbreak{} i\allowbreak{} d\allowbreak{}i\allowbreak{}v\allowbreak{} 2\allowbreak{})}. \\
\texttt{read\_\allowbreak{}eval} & 52-57 & test fixture & The law for the three slices. \\
\texttt{layout} & 60-63 & test fixture & The toy's layout. \\
\texttt{constraint} & 66-66 & test fixture & Column 2 is Boolean: \texttt{X₂² − X₂}. \\
\texttt{flush} & 77-78 & test fixture & The one flush: \texttt{(\allowbreak{}X₀,\allowbreak{} 0,\allowbreak{} …)}, pushed. \\
\texttt{boundary} & 89-92 & test fixture & The one boundary block: pulls \texttt{(\allowbreak{}[\allowbreak{}1\allowbreak{},\allowbreak{} 1\allowbreak{}]\allowbreak{},\allowbreak{} 0\allowbreak{},\allowbreak{} …\allowbreak{})}. \\
\texttt{toy} & 95-111 & test fixture & The toy instance: \texttt{d = 2}, column 1 a count column, the public line on column 2 with cells \texttt{(\allowbreak{}input,\allowbreak{} 0)}, \texttt{a\allowbreak{}u\allowbreak{}x\allowbreak{} :\allowbreak{}=\allowbreak{} T\allowbreak{}r\allowbreak{}u\allowbreak{}e}. \\
\texttt{honest} & 114-114 & test fixture & The honest stack \texttt{[\allowbreak{}1\allowbreak{},\allowbreak{} 1\allowbreak{},\allowbreak{} 1\allowbreak{},\allowbreak{} 1\allowbreak{},\allowbreak{} 1\allowbreak{},\allowbreak{} 0\allowbreak{},\allowbreak{} 0\allowbreak{},\allowbreak{} 0\allowbreak{}]}. \\
\end{longtable}}

\subsection{\texttt{LeanerVM/Protocol/ToArkLib/Oracles.lean}}\label{sec:gen-cat-spine-Oracles}

\par\medskip\noindent\textbf{\lean{NoOracle}}\hfill\textsf{\small [load-bearing]}\par\nopagebreak

\begin{leancode}
abbrev NoOracle : Fin 0 → Type := fun i ↦ i.elim0
\end{leancode}

\src{LeanerVM/Protocol/ToArkLib/Oracles.lean:26-26 at b435631}

The oracle family with no oracle (before the commit). \emph{Depends on:} —.

\par\medskip\noindent\textbf{\lean{OneOracle}}\hfill\textsf{\small [load-bearing]}\par\nopagebreak

\begin{leancode}
abbrev OneOracle (M : Type) : Fin 1 → Type := fun _ ↦ M
\end{leancode}

\src{LeanerVM/Protocol/ToArkLib/Oracles.lean:29-29 at b435631}

The oracle family with exactly one oracle of type \texttt{M} (the stack, from the commit on). \emph{Depends on:} —.

\par\medskip\noindent\textbf{Helpers and test fixtures of \file{Oracles.lean}.}

{\footnotesize
\begin{longtable}{@{}>{\raggedright\arraybackslash}p{0.300\linewidth}>{\raggedright\arraybackslash}p{0.100\linewidth}>{\raggedright\arraybackslash}p{0.160\linewidth}>{\raggedright\arraybackslash}p{0.360\linewidth}@{}}
\toprule
Declaration & Lines & Class & Meaning \\
\midrule
\endfirsthead
\toprule
Declaration & Lines & Class & Meaning \\
\midrule
\endhead
\bottomrule
\endfoot
\texttt{n\allowbreak{}o\allowbreak{}O\allowbreak{}r\allowbreak{}a\allowbreak{}c\allowbreak{}l\allowbreak{}e\allowbreak{}\_\allowbreak{}e\allowbreak{}q} & 32-32 & helper & Any two empty oracle families are equal. \\
\end{longtable}}

\subsection{\texttt{LeanerVM/Protocol/ToArkLib/Component.lean}}\label{sec:gen-cat-spine-Component}

\par\medskip\noindent\textbf{\lean{Component.Def}}\hfill\textsf{\small [interface]}\par\nopagebreak

\begin{leancode}
structure Def (StmtIn : Type) {ιi : Type} (OStmtIn : ιi → Type) (WitIn : Type)
    (StmtOut : Type) {ιo : Type} (OStmtOut : ιo → Type) (WitOut : Type)
    [∀ i, OracleInterface (OStmtIn i)] [∀ i, OracleInterface (OStmtOut i)] where
  n : ℕ
  pSpec : ProtocolSpec n
  [msgOracle : ∀ i, OracleInterface (pSpec.Message i)]
  [chalSample : ∀ i, SampleableType (pSpec.Challenge i)]
  red : OracleReduction []ₒ StmtIn OStmtIn WitIn StmtOut OStmtOut WitOut pSpec
  err : pSpec.ChallengeIdx → ℝ≥0
\end{leancode}

\src{LeanerVM/Protocol/ToArkLib/Component.lean:52-66 at b435631}

A component: its number of rounds, its schedule (who speaks, what type), the instances that every message can be queried and every challenge sampled, the honest prover with the verifier (an ArkLib \texttt{OracleReduction} over the empty shared oracle), and the knowledge error it declares per challenge. \emph{Depends on:} ArkLib \texttt{ProtocolSpec}, \texttt{OracleReduction}, \texttt{OracleInterface}, VCVio \texttt{SampleableType}.

\par\medskip\noindent\textbf{\lean{Component.Complete}}\hfill\textsf{\small [interface]}\par\nopagebreak

\begin{leancode}
structure Complete (D : Def StmtIn OStmtIn WitIn StmtOut OStmtOut WitOut)
    (relIn : Set ((StmtIn × ∀ i, OStmtIn i) × WitIn))
    (relOut : Set ((StmtOut × ∀ i, OStmtOut i) × WitOut)) where
  outputPure : D.red.prover.OutputIsPure
  guarded : D.red.toReduction.verifier.GuardedForm
  complete : ∀ {σ : Type} (init : ProbComp σ) (impl : QueryImpl []ₒ (StateT σ ProbComp)),
    D.red.perfectCompleteness init impl relIn relOut
\end{leancode}

\src{LeanerVM/Protocol/ToArkLib/Component.lean:76-85 at b435631}

The completeness half a phase must supply: the prover's output is pure (redundant, D.7), the verifier is a check followed by a verdict (data), and perfect completeness from any shared state. \emph{Depends on:} \texttt{Def}, ArkLib \texttt{OutputIsPure}, \texttt{GuardedForm}, \texttt{perfectCompleteness}.

\par\medskip\noindent\textbf{\lean{Component.Security}}\hfill\textsf{\small [interface]}\par\nopagebreak

\begin{leancode}
structure Security (D : Def StmtIn OStmtIn WitIn StmtOut OStmtOut WitOut)
    (relIn : Set ((StmtIn × ∀ i, OStmtIn i) × WitIn))
    (relOut : Set ((StmtOut × ∀ i, OStmtOut i) × WitOut))
    extends Complete D relIn relOut where
  witMid : Fin (D.n + 1) → Type
  extractor : Extractor.RoundByRound (OracleSpec.emptySpec.{0, 0}) (StmtIn × ∀ i, OStmtIn i)
    WitIn WitOut D.pSpec witMid
  kSF : ∀ {σ : Type} (init : ProbComp σ) (impl : QueryImpl []ₒ (StateT σ ProbComp)),
    D.red.verifier.toVerifier.KnowledgeStateFunction init impl relIn relOut extractor
  rbr : ∀ {σ : Type} (init : ProbComp σ) (impl : QueryImpl []ₒ (StateT σ ProbComp)),
    D.red.verifier.toVerifier.rbrKnowledgeSoundnessWorstCaseWith init impl relIn relOut witMid
      extractor (kSF init impl) D.err
\end{leancode}

\src{LeanerVM/Protocol/ToArkLib/Component.lean:90-106 at b435631}

The security half: completeness, the intermediate witness types, a round-by-round extractor, its knowledge state function from any shared state, and worst-case round-by-round knowledge soundness for them at the declared error. \emph{Depends on:} \texttt{Complete}, ArkLib \texttt{Extractor.\allowbreak{}RoundByRound}, \texttt{KnowledgeStateFunction}, \texttt{r\allowbreak{}b\allowbreak{}r\allowbreak{}K\allowbreak{}n\allowbreak{}o\allowbreak{}w\allowbreak{}l\allowbreak{}e\allowbreak{}d\allowbreak{}g\allowbreak{}e\allowbreak{}S\allowbreak{}o\allowbreak{}u\allowbreak{}n\allowbreak{}d\allowbreak{}n\allowbreak{}e\allowbreak{}s\allowbreak{}s\allowbreak{}W\allowbreak{}o\allowbreak{}r\allowbreak{}s\allowbreak{}t\allowbreak{}C\allowbreak{}a\allowbreak{}s\allowbreak{}e\allowbreak{}W\allowbreak{}i\allowbreak{}t\allowbreak{}h}.

\par\medskip\noindent\textbf{\lean{Component.Def.append}}\hfill\textsf{\small [load-bearing]}\par\nopagebreak

\begin{leancode}
def Def.append (D₁ : Def Stmt₁ OStmt₁ Wit₁ Stmt₂ OStmt₂ Wit₂)
    (D₂ : Def Stmt₂ OStmt₂ Wit₂ Stmt₃ OStmt₃ Wit₃) : Def Stmt₁ OStmt₁ Wit₁ Stmt₃ OStmt₃ Wit₃ where
  n := D₁.n + D₂.n
  pSpec := D₁.pSpec ++ₚ D₂.pSpec
  red := D₁.red.append D₂.red
  err := Sum.elim D₁.err D₂.err ∘ ChallengeIdx.sumEquiv.symm
\end{leancode}

\src{LeanerVM/Protocol/ToArkLib/Component.lean:143-148 at b435631}

Two components in sequence: rounds add, schedules concatenate, reductions append (ArkLib), each challenge keeps its component's error. \emph{Depends on:} ArkLib \texttt{OracleReduction.\allowbreak{}append}, \texttt{ChallengeIdx.\allowbreak{}sumEquiv}.

\par\medskip\noindent\textbf{Helpers and test fixtures of \file{Component.lean}.}

{\footnotesize
\begin{longtable}{@{}>{\raggedright\arraybackslash}p{0.300\linewidth}>{\raggedright\arraybackslash}p{0.100\linewidth}>{\raggedright\arraybackslash}p{0.160\linewidth}>{\raggedright\arraybackslash}p{0.360\linewidth}@{}}
\toprule
Declaration & Lines & Class & Meaning \\
\midrule
\endfirsthead
\toprule
Declaration & Lines & Class & Meaning \\
\midrule
\endhead
\bottomrule
\endfoot
\texttt{C\allowbreak{}o\allowbreak{}m\allowbreak{}p\allowbreak{}o\allowbreak{}n\allowbreak{}e\allowbreak{}n\allowbreak{}t\allowbreak{}.\allowbreak{}s\allowbreak{}t\allowbreak{}a\allowbreak{}t\allowbreak{}e\allowbreak{}F\allowbreak{}u\allowbreak{}n\allowbreak{}c\allowbreak{}t\allowbreak{}i\allowbreak{}o\allowbreak{}n\allowbreak{}O\allowbreak{}f\allowbreak{}E\allowbreak{}q} & 110-120 & helper (in the named statement) & Transport a knowledge state function along an equality of verifiers (needed because ArkLib's appended oracle verifier is only propositionally the appended ordinary verifier). \\
\texttt{r\allowbreak{}b\allowbreak{}r\allowbreak{}K\allowbreak{}n\allowbreak{}o\allowbreak{}w\allowbreak{}l\allowbreak{}e\allowbreak{}d\allowbreak{}g\allowbreak{}e\allowbreak{}S\allowbreak{}o\allowbreak{}u\allowbreak{}n\allowbreak{}d\allowbreak{}n\allowbreak{}e\allowbreak{}s\allowbreak{}s\allowbreak{}W\allowbreak{}o\allowbreak{}r\allowbreak{}s\allowbreak{}t\allowbreak{}C\allowbreak{}a\allowbreak{}s\allowbreak{}e\allowbreak{}W\allowbreak{}i\allowbreak{}t\allowbreak{}h\allowbreak{}\_\allowbreak{}o\allowbreak{}f\allowbreak{}\_\allowbreak{}e\allowbreak{}q} & 124-135 & helper & The same transport for the bound. \\
\texttt{C\allowbreak{}o\allowbreak{}m\allowbreak{}p\allowbreak{}o\allowbreak{}n\allowbreak{}e\allowbreak{}n\allowbreak{}t\allowbreak{}.\allowbreak{}g\allowbreak{}u\allowbreak{}a\allowbreak{}r\allowbreak{}d\allowbreak{}e\allowbreak{}d\allowbreak{}A\allowbreak{}p\allowbreak{}p\allowbreak{}e\allowbreak{}n\allowbreak{}d} & 156-159 & helper (in the named statement) & The guarded form of an appended verifier from those of its parts (ArkLib's \texttt{G\allowbreak{}u\allowbreak{}a\allowbreak{}r\allowbreak{}d\allowbreak{}e\allowbreak{}d\allowbreak{}F\allowbreak{}o\allowbreak{}r\allowbreak{}m\allowbreak{}.\allowbreak{}a\allowbreak{}p\allowbreak{}p\allowbreak{}e\allowbreak{}n\allowbreak{}d}, cast along the append equation). \\
\texttt{C\allowbreak{}o\allowbreak{}m\allowbreak{}p\allowbreak{}o\allowbreak{}n\allowbreak{}e\allowbreak{}n\allowbreak{}t\allowbreak{}.\allowbreak{}C\allowbreak{}o\allowbreak{}m\allowbreak{}p\allowbreak{}l\allowbreak{}e\allowbreak{}t\allowbreak{}e\allowbreak{}.\allowbreak{}a\allowbreak{}p\allowbreak{}p\allowbreak{}e\allowbreak{}n\allowbreak{}d} & 162-169 & helper (in the named statement) & Completeness composes: purity, guarded form, and ArkLib's guarded-append completeness theorem. \\
\texttt{C\allowbreak{}o\allowbreak{}m\allowbreak{}p\allowbreak{}o\allowbreak{}n\allowbreak{}e\allowbreak{}n\allowbreak{}t\allowbreak{}.\allowbreak{}S\allowbreak{}e\allowbreak{}c\allowbreak{}u\allowbreak{}r\allowbreak{}i\allowbreak{}t\allowbreak{}y\allowbreak{}.\allowbreak{}a\allowbreak{}p\allowbreak{}p\allowbreak{}e\allowbreak{}n\allowbreak{}d} & 174-184 & helper (in the named statement) & Security composes: extractors appended through the first verdict (ArkLib), state functions by the ported \texttt{a\allowbreak{}p\allowbreak{}p\allowbreak{}e\allowbreak{}n\allowbreak{}d\allowbreak{}G\allowbreak{}u\allowbreak{}a\allowbreak{}r\allowbreak{}d\allowbreak{}e\allowbreak{}d}, the bound by the ported theorem. \\
\end{longtable}}

\subsection{\texttt{LeanerVM/Protocol/ToArkLib/KnowledgeAppend.lean}}\label{sec:gen-cat-spine-KnowledgeAppend}

\par\medskip\noindent\textbf{Helpers and test fixtures of \file{KnowledgeAppend.lean}.}

{\footnotesize
\begin{longtable}{@{}>{\raggedright\arraybackslash}p{0.300\linewidth}>{\raggedright\arraybackslash}p{0.100\linewidth}>{\raggedright\arraybackslash}p{0.160\linewidth}>{\raggedright\arraybackslash}p{0.360\linewidth}@{}}
\toprule
Declaration & Lines & Class & Meaning \\
\midrule
\endfirsthead
\toprule
Declaration & Lines & Class & Meaning \\
\midrule
\endhead
\bottomrule
\endfoot
\texttt{V\allowbreak{}e\allowbreak{}r\allowbreak{}i\allowbreak{}f\allowbreak{}i\allowbreak{}e\allowbreak{}r\allowbreak{}.\allowbreak{}K\allowbreak{}n\allowbreak{}o\allowbreak{}w\allowbreak{}l\allowbreak{}e\allowbreak{}d\allowbreak{}g\allowbreak{}e\allowbreak{}S\allowbreak{}t\allowbreak{}a\allowbreak{}t\allowbreak{}e\allowbreak{}F\allowbreak{}u\allowbreak{}n\allowbreak{}c\allowbreak{}t\allowbreak{}i\allowbreak{}o\allowbreak{}n\allowbreak{}.\allowbreak{}a\allowbreak{}p\allowbreak{}p\allowbreak{}e\allowbreak{}n\allowbreak{}d\allowbreak{}G\allowbreak{}u\allowbreak{}a\allowbreak{}r\allowbreak{}d\allowbreak{}e\allowbreak{}d} & 474-488 & helper (in the named statement) & The knowledge state function of two appended verifiers, the first guarded: the first's state up to the seam, then the first check together with the second's state on the first verdict (ported from ArkLib \#615). \\
\texttt{V\allowbreak{}e\allowbreak{}r\allowbreak{}i\allowbreak{}f\allowbreak{}i\allowbreak{}e\allowbreak{}r\allowbreak{}.\allowbreak{}a\allowbreak{}p\allowbreak{}p\allowbreak{}e\allowbreak{}n\allowbreak{}d\allowbreak{}\_\allowbreak{}r\allowbreak{}b\allowbreak{}r\allowbreak{}K\allowbreak{}n\allowbreak{}o\allowbreak{}w\allowbreak{}l\allowbreak{}e\allowbreak{}d\allowbreak{}g\allowbreak{}e\allowbreak{}S\allowbreak{}o\allowbreak{}u\allowbreak{}n\allowbreak{}d\allowbreak{}n\allowbreak{}e\allowbreak{}s\allowbreak{}s\allowbreak{}W\allowbreak{}o\allowbreak{}r\allowbreak{}s\allowbreak{}t\allowbreak{}C\allowbreak{}a\allowbreak{}s\allowbreak{}e\allowbreak{}W\allowbreak{}i\allowbreak{}t\allowbreak{}h\allowbreak{}\_\allowbreak{}o\allowbreak{}f\allowbreak{}\_\allowbreak{}g\allowbreak{}u\allowbreak{}a\allowbreak{}r\allowbreak{}d\allowbreak{}e\allowbreak{}d\allowbreak{}\_\allowbreak{}f\allowbreak{}i\allowbreak{}r\allowbreak{}s\allowbreak{}t} & 510-582 & helper & Round-by-round knowledge soundness composes across an append whose first verifier is guarded, each challenge keeping its error (ported from ArkLib \#615; admitted upstream at the pin). \\
\end{longtable}}

\subsection{\texttt{LeanerVM/Protocol/ToArkLib/PassThrough.lean}}\label{sec:gen-cat-spine-PassThrough}

\par\medskip\noindent\textbf{\lean{Component.passThrough}}\hfill\textsf{\small [interface]}\par\nopagebreak

\begin{leancode}
def passThrough (f : StmtIn → StmtOut) : Def StmtIn OStmt W StmtOut OStmt W where
  n := 0
  pSpec := !p[]
  red := ⟨passThroughProver OStmt f, passThroughVerifier OStmt f⟩
  err := fun i ↦ Fin.elim0 i.1
\end{leancode}

\src{LeanerVM/Protocol/ToArkLib/PassThrough.lean:53-57 at b435631}

The component with no round that maps the statement and keeps the oracles and the witness: the shape of bookkeeping steps. \emph{Depends on:} \texttt{passThroughProver}, \texttt{passThroughVerifier}, \texttt{keepOracles}.

\par\medskip\noindent\textbf{\lean{Component.passThroughComplete}}\hfill\textsf{\small [interface]}\par\nopagebreak

\begin{leancode}
def passThroughComplete (f : StmtIn → StmtOut)
    {relIn : Set ((StmtIn × ∀ i, OStmt i) × W)} {relOut : Set ((StmtOut × ∀ i, OStmt i) × W)}
    (h : ∀ s o w, ((s, o), w) ∈ relIn → ((f s, o), w) ∈ relOut) :
    Complete (passThrough OStmt f) relIn relOut where
  outputPure := ⟨_, fun _ ↦ rfl⟩
  guarded := (passThroughPure OStmt f).toGuardedForm
\end{leancode}

\src{LeanerVM/Protocol/ToArkLib/PassThrough.lean:81-95 at b435631}

Its completeness, whenever the map carries the input relation into the output relation. \emph{Depends on:} \texttt{passThrough}, \texttt{GuardedVerdict}.

\par\medskip\noindent\textbf{\lean{Component.passThroughSecurity}}\hfill\textsf{\small [interface]}\par\nopagebreak

\begin{leancode}
def passThroughSecurity (hc : ∀ s o w, ((s, o), w) ∈ relIn → ((f s, o), w) ∈ relOut) :
    Security (passThrough OStmt f) relIn relOut where
  toComplete := passThroughComplete OStmt f hc
  witMid := fun _ ↦ W
  extractor := passThroughExtractor OStmt
  kSF := passThroughStateFunction OStmt f h
\end{leancode}

\src{LeanerVM/Protocol/ToArkLib/PassThrough.lean:132-138 at b435631}

Its security at error zero, with the extractor that keeps the witness, whenever the map also reflects the output relation into the input relation. \emph{Depends on:} \texttt{passThroughExtractor}, \texttt{passThroughStateFunction}, \texttt{passThrough\_\allowbreak{}rbr}.

\subsection{\texttt{LeanerVM/Protocol/ToArkLib/SendOracle.lean}}\label{sec:gen-cat-spine-SendOracle}

\par\medskip\noindent\textbf{\lean{Component.sendOracle}}\hfill\textsf{\small [load-bearing]}\par\nopagebreak

\begin{leancode}
def sendOracle : Def S NoOracle M S (OneOracle M) Unit where
  n := 1
  pSpec := sendSpec M
  red := ⟨sendProver S M, sendVerifier S M⟩
  err := fun _ ↦ 0
\end{leancode}

\src{LeanerVM/Protocol/ToArkLib/SendOracle.lean:61-65 at b435631}

The component whose one message becomes the one oracle: the prover sends the witness, the verifier keeps the statement and exposes the message; no challenge, error zero. \emph{Depends on:} \texttt{sendSpec}, \texttt{sendProver}, \texttt{sendVerifier}, \texttt{sendEmbedding}.

\par\medskip\noindent\textbf{\lean{Component.sendOracle_relOut}}\hfill\textsf{\small [load-bearing (named form)]}\par\nopagebreak

\begin{leancode}
def sendOracle_relOut (rel : Set ((S × ∀ i, NoOracle i) × M)) :
    Set ((S × ∀ i, OneOracle M i) × Unit) :=
  {p | ((p.1.1, fun i ↦ i.elim0), p.1.2 0) ∈ rel}
\end{leancode}

\src{LeanerVM/Protocol/ToArkLib/SendOracle.lean:71-73 at b435631}

The input relation read on the oracle instead of the witness (definitionally \texttt{Seam.\allowbreak{}commit} at \texttt{M3Rel}). \emph{Depends on:} —.

\par\medskip\noindent\textbf{\lean{Component.sendOracleComplete}}\hfill\textsf{\small [load-bearing (named form)]}\par\nopagebreak

\begin{leancode}
def sendOracleComplete (rel : Set ((S × ∀ i, NoOracle i) × M)) :
    Complete (sendOracle S M) rel (sendOracle_relOut rel) where
  outputPure := sendOracle_outputPure
  guarded := (sendVerifierPure (S := S) (M := M)).toGuardedForm
\end{leancode}

\src{LeanerVM/Protocol/ToArkLib/SendOracle.lean:121-125 at b435631}

Its completeness half. \emph{Depends on:} \texttt{sendVerifierPure}, \texttt{sendOracle\_\allowbreak{}complete}.

\par\medskip\noindent\textbf{\lean{Component.sendExtractor}}\hfill\textsf{\small [load-bearing (named form)]}\par\nopagebreak

\begin{leancode}
def sendExtractor : Extractor.RoundByRound (OracleSpec.emptySpec.{0, 0}) (S × ∀ i, NoOracle i) M
    Unit (sendSpec M) (sendWitMid M) where
  extractMid := fun m _ tr _ ↦ tr ⟨0, Nat.succ_pos m.val⟩
  extractOut := fun _ tr _ ↦ tr 0
\end{leancode}

\src{LeanerVM/Protocol/ToArkLib/SendOracle.lean:134-138 at b435631}

Its extractor: read the message, at every round. \emph{Depends on:} ArkLib \texttt{Extractor.\allowbreak{}RoundByRound}.

\par\medskip\noindent\textbf{\lean{Component.sendStateFunction}}\hfill\textsf{\small [load-bearing (named form)]}\par\nopagebreak

\begin{leancode}
def sendStateFunction :
    (sendVerifier S M).toVerifier.KnowledgeStateFunction init impl rel (sendOracle_relOut rel)
      (sendExtractor (S := S) (M := M)) where
  toFun := fun m stmt tr w ↦
    if h : m.val = 0 then (stmt, w) ∈ rel else (stmt, tr ⟨0, Nat.pos_of_ne_zero h⟩) ∈ rel
\end{leancode}

\src{LeanerVM/Protocol/ToArkLib/SendOracle.lean:145-161 at b435631}

Its knowledge state function: before the message, the relation of the candidate witness; after it, the relation of the message. \emph{Depends on:} \texttt{sendExtractor}, \texttt{GuardedVerdict}.

\par\medskip\noindent\textbf{\lean{Component.sendOracleSecurity}}\hfill\textsf{\small [load-bearing (named form)]}\par\nopagebreak

\begin{leancode}
def sendOracleSecurity : Security (sendOracle S M) rel (sendOracle_relOut rel) where
  toComplete := sendOracleComplete rel
  witMid := sendWitMid M
  extractor := sendExtractor
  kSF := sendStateFunction rel
\end{leancode}

\src{LeanerVM/Protocol/ToArkLib/SendOracle.lean:172-177 at b435631}

Its security half at error zero. \emph{Depends on:} \texttt{sendOracleComplete}, \texttt{sendExtractor}, \texttt{sendStateFunction}, \texttt{sendOracle\_\allowbreak{}rbr}.

\subsection{\texttt{LeanerVM/Protocol/ToArkLib/Refinement.lean}}\label{sec:gen-cat-spine-Refinement}

\par\medskip\noindent\textbf{\lean{Refinement}}\hfill\textsf{\small [interface (the adaptor)]}\par\nopagebreak

\begin{leancode}
structure Refinement {Stmt : Type} {W₁ W₂ : Type _} (R : Set (Stmt × W₁))
    (S : Set (Stmt × W₂)) where
  map : Stmt → W₁ → W₂
  map_valid : ∀ x w, (x, w) ∈ R → (x, map x w) ∈ S
\end{leancode}

\src{LeanerVM/Protocol/ToArkLib/Refinement.lean:31-36 at b435631}

A map on witnesses between two relations over the same statements, sending valid witnesses to valid witnesses; the two witness types may live in different universes (D.7). \emph{Depends on:} —.

\par\medskip\noindent\textbf{\lean{Refinement.map_option_valid}}\hfill\textsf{\small [interface (the adaptor)]}\par\nopagebreak

\begin{leancode}
theorem map_option_valid {R : Set (Stmt × W₁)} {S : Set (Stmt × W₂)} (f : Refinement R S)
    (x : Stmt) (w? : Option W₁) (h : ∀ w ∈ w?, (x, w) ∈ R) :
    ∀ w' ∈ w?.map (f.map x), (x, w') ∈ S :=
\end{leancode}

\src{LeanerVM/Protocol/ToArkLib/Refinement.lean:55-60 at b435631}

On an extractor's output slot: if every witness the slot may hold is valid for \texttt{R}, every witness of the mapped slot is valid for \texttt{S}. Pointwise; the probabilistic form is D.7. \emph{Depends on:} \texttt{Refinement}.

\par\medskip\noindent\textbf{Helpers and test fixtures of \file{Refinement.lean}.}

{\footnotesize
\begin{longtable}{@{}>{\raggedright\arraybackslash}p{0.300\linewidth}>{\raggedright\arraybackslash}p{0.100\linewidth}>{\raggedright\arraybackslash}p{0.160\linewidth}>{\raggedright\arraybackslash}p{0.360\linewidth}@{}}
\toprule
Declaration & Lines & Class & Meaning \\
\midrule
\endfirsthead
\toprule
Declaration & Lines & Class & Meaning \\
\midrule
\endhead
\bottomrule
\endfoot
\texttt{E\allowbreak{}x\allowbreak{}t\allowbreak{}r\allowbreak{}a\allowbreak{}c\allowbreak{}t\allowbreak{}o\allowbreak{}r\allowbreak{}.\allowbreak{}S\allowbreak{}t\allowbreak{}r\allowbreak{}a\allowbreak{}i\allowbreak{}g\allowbreak{}h\allowbreak{}t\allowbreak{}l\allowbreak{}i\allowbreak{}n\allowbreak{}e\allowbreak{}.\allowbreak{}m\allowbreak{}a\allowbreak{}p} & 68-71 & helper (no consumer) & Post-compose ArkLib's straight-line extractor with a witness map (same universe). \\
\end{longtable}}

\subsection{\texttt{LeanerVM/Protocol/ToArkLib/GuardedVerdict.lean}}\label{sec:gen-cat-spine-GuardedVerdict}

\par\medskip\noindent\textbf{Helpers and test fixtures of \file{GuardedVerdict.lean}.}

{\footnotesize
\begin{longtable}{@{}>{\raggedright\arraybackslash}p{0.300\linewidth}>{\raggedright\arraybackslash}p{0.100\linewidth}>{\raggedright\arraybackslash}p{0.160\linewidth}>{\raggedright\arraybackslash}p{0.360\linewidth}@{}}
\toprule
Declaration & Lines & Class & Meaning \\
\midrule
\endfirsthead
\toprule
Declaration & Lines & Class & Meaning \\
\midrule
\endhead
\bottomrule
\endfoot
\texttt{V\allowbreak{}e\allowbreak{}r\allowbreak{}i\allowbreak{}f\allowbreak{}i\allowbreak{}e\allowbreak{}r\allowbreak{}.\allowbreak{}G\allowbreak{}u\allowbreak{}a\allowbreak{}r\allowbreak{}d\allowbreak{}e\allowbreak{}d\allowbreak{}F\allowbreak{}o\allowbreak{}r\allowbreak{}m\allowbreak{}.\allowbreak{}o\allowbreak{}f\allowbreak{}\_\allowbreak{}p\allowbreak{}r\allowbreak{}o\allowbreak{}b\allowbreak{}E\allowbreak{}v\allowbreak{}e\allowbreak{}n\allowbreak{}t\allowbreak{}\_\allowbreak{}p\allowbreak{}o\allowbreak{}s} & 40-69 & helper & If a guarded verifier can output a statement satisfying \texttt{P}, its check passes and its verdict satisfies \texttt{P}: the last obligation of every knowledge state function. \\
\texttt{R\allowbreak{}e\allowbreak{}d\allowbreak{}u\allowbreak{}c\allowbreak{}t\allowbreak{}i\allowbreak{}o\allowbreak{}n\allowbreak{}.\allowbreak{}m\allowbreak{}e\allowbreak{}m\allowbreak{}\_\allowbreak{}s\allowbreak{}u\allowbreak{}p\allowbreak{}p\allowbreak{}o\allowbreak{}r\allowbreak{}t\allowbreak{}\_\allowbreak{}r\allowbreak{}u\allowbreak{}n\allowbreak{}\_\allowbreak{}o\allowbreak{}f\allowbreak{}\_\allowbreak{}g\allowbreak{}u\allowbreak{}a\allowbreak{}r\allowbreak{}d\allowbreak{}e\allowbreak{}d} & 73-112 & helper & Every outcome of a run with a guarded verifier is a prover run with the verdict where the check passes, a rejection otherwise: how perfect completeness is proved. \\
\end{longtable}}

\subsection{\texttt{LeanerVM/Protocol/ToArkLib/KeepOracles.lean}}\label{sec:gen-cat-spine-KeepOracles}

\par\medskip\noindent\textbf{Helpers and test fixtures of \file{KeepOracles.lean}.}

{\footnotesize
\begin{longtable}{@{}>{\raggedright\arraybackslash}p{0.300\linewidth}>{\raggedright\arraybackslash}p{0.100\linewidth}>{\raggedright\arraybackslash}p{0.160\linewidth}>{\raggedright\arraybackslash}p{0.360\linewidth}@{}}
\toprule
Declaration & Lines & Class & Meaning \\
\midrule
\endfirsthead
\toprule
Declaration & Lines & Class & Meaning \\
\midrule
\endhead
\bottomrule
\endfoot
\texttt{k\allowbreak{}e\allowbreak{}e\allowbreak{}p\allowbreak{}O\allowbreak{}r\allowbreak{}a\allowbreak{}c\allowbreak{}l\allowbreak{}e\allowbreak{}s} & 34-37 & helper & The output-oracle description \enquote{the output oracles are the input oracles}. \\
\texttt{O\allowbreak{}r\allowbreak{}a\allowbreak{}c\allowbreak{}l\allowbreak{}e\allowbreak{}V\allowbreak{}e\allowbreak{}r\allowbreak{}i\allowbreak{}f\allowbreak{}i\allowbreak{}e\allowbreak{}r\allowbreak{}.\allowbreak{}m\allowbreak{}a\allowbreak{}t\allowbreak{}e\allowbreak{}r\allowbreak{}i\allowbreak{}a\allowbreak{}l\allowbreak{}i\allowbreak{}z\allowbreak{}e\allowbreak{}O\allowbreak{}u\allowbreak{}t\allowbreak{}p\allowbreak{}u\allowbreak{}t\allowbreak{}\_\allowbreak{}o\allowbreak{}f\allowbreak{}\_\allowbreak{}k\allowbreak{}e\allowbreak{}e\allowbreak{}p\allowbreak{}O\allowbreak{}r\allowbreak{}a\allowbreak{}c\allowbreak{}l\allowbreak{}e\allowbreak{}s} & 42-50 & helper & Such a verifier outputs the oracles it was given. \\
\end{longtable}}

\subsection{\texttt{LeanerVM/Protocol/Field.lean}}\label{sec:gen-cat-spine-Field}

\par\medskip\noindent\textbf{\lean{Column}}\hfill\textsf{\small [load-bearing]}\par\nopagebreak

\begin{leancode}
structure Column (n : ℕ) where
  values : CMlPolynomialEval K n
\end{leancode}

\src{LeanerVM/Protocol/Field.lean:67-69 at b435631}

A table of \texttt{2\textasciicircum{}n} values of \texttt{K}, indexed by the cube \texttt{\{0,\allowbreak{}1\}\textasciicircum{}n} low bit first: what is committed. Wrapped in a structure so that its oracle interface is the evaluation one, not ArkLib's position query on \texttt{Vector}. \emph{Depends on:} CompPoly \texttt{CMlPolynomialEval}.

\par\medskip\noindent\textbf{\lean{evalOracle}}\hfill\textsf{\small [load-bearing]}\par\nopagebreak

\begin{leancode}
instance evalOracle (n : ℕ) : OracleInterface (Column n) where
  Query := Vector E n
  toOC :=
    { spec := (Vector E n) →ₒ E
      impl := fun r ↦ do return CMlPolynomialEval.eval₂Mle (← read).values (algebraMap K E) r }
\end{leancode}

\src{LeanerVM/Protocol/Field.lean:102-106 at b435631}

How the verifier queries a committed column: a query is a point \texttt{r ∈ E\textasciicircum{}n}, the answer is the multilinear extension \texttt{q̃(\allowbreak{}r)} with the entries lifted from \texttt{K} to \texttt{E}. Every seam is read through it. \emph{Depends on:} \texttt{Column}, CompPoly \texttt{eval₂Mle}, ArkLib \texttt{OracleInterface}.

\par\medskip\noindent\textbf{\lean{instSampleableTypeE}}\hfill\textsf{\small [load-bearing]}\par\nopagebreak

\begin{leancode}
instance instSampleableTypeE : SampleableType E := SampleableType.ofEquiv limbsEquiv
\end{leancode}

\src{LeanerVM/Protocol/Field.lean:93-93 at b435631}

Uniform sampling of a challenge in \texttt{E} as three independent uniform limbs (VCVio needs a sampler for every challenge type). \emph{Depends on:} \texttt{limbsEquiv}, VCVio \texttt{SampleableType}.
